% !TeX root = ../main.tex
%%% ---------------------------------------------------------------------------
%%% 思考题
%%%
%%% 大多数课程的实验报告都有思考题。这一节是拉分项：
%%% 助教看的是你有没有把原理真正想明白，答案要有推导或数据支撑，
%%% 不要写成「我觉得……」。
%%%
%%% type = tech（技术报告）时通常没有这一节，把 main.tex 里对应的
%%% \input 注释掉即可。
%%% ---------------------------------------------------------------------------
\section{思考题}
\label{sec:questions}

\subsection*{1. 为什么不能用测试集选择最优 epoch？}

用测试集选点相当于让测试集参与了模型选择，此时测试准确率不再是泛化性能的
无偏估计。设候选模型有 $k$ 个，即使它们泛化能力相同，取测试集上的最大值也会
系统性地高于真实期望，偏差量随 $k$ 增大而增大。本实验中每次运行有
\num{100} 个 epoch 候选，若用测试集选点，估计的乐观偏差可达零点几个百分点，
与\cref{tab:summary} 中优化器之间的真实差异同量级，足以颠倒结论。

\subsection*{2. 为什么 Adam 和 AdamW 的权重衰减取值不同却仍可比较？}

因为二者的 $\lambda$ 作用在不同的位置。Adam 把 $\lambda \params$ 加进梯度，
随后被二阶矩 $\sqrt{\hat{\vect{v}}_t}$ 整体缩放，实际衰减强度是
$\lambda / (\sqrt{\hat{\vect{v}}_t} + \epsilon)$，随参数而异；AdamW 按
\cref{eq:adamw} 直接作用于参数，强度恒为 $\lr\lambda$。因此两者的
$\lambda$ 不是同一个物理量，可比较的是各自调优后的最好结果，
而不是相同 $\lambda$ 下的结果。

\subsection*{3. 若把重复次数从 5 增加到 20，不确定度会如何变化？}

由\cref{eq:uncertainty}，$u(\bar{x}) = s / \sqrt{n}$，在 $s$ 不变的前提下

\begin{equation}
  \label{eq:u-scaling}
  \frac{u_{20}}{u_{5}} = \sqrt{\frac{5}{20}} = \frac{1}{2},
\end{equation}

即不确定度减半，从约 \qty{0.07}{\percent} 降到 \qty{0.035}{\percent}。
但计算量增至 4 倍，而系统误差（\cref{tab:error} 中未消除的两项）不会随
$n$ 减小。因此当随机不确定度已经小于系统误差时，继续增加重复次数收益很低，
应当转而去消除系统误差。

\subsection*{4. 实验中哪一步最可能被同学做错，为什么？}

固定随机种子这一步（\cref{sec:procedure} 第 1 步）。多数人只设了
\lstinline|torch.manual_seed|，忘了 \lstinline|numpy| 与 DataLoader 的
\lstinline|worker_init_fn|，导致三种优化器实际拿到的数据顺序不同。
这样测出的差异里混入了数据顺序的贡献，而这一项在本实验中的量级
（约 \qty{0.15}{\percent}，即\cref{eq:sgdm-std} 的 $s$）恰好与
SGD-M 和 AdamW 的真实差异相当，足以让结论失真。
